%---------- Inleiding ---------------------------------------------------------

% TODO: Is dit voorstel gebaseerd op een paper van Research Methods die je
% vorig jaar hebt ingediend? Heb je daarbij eventueel samengewerkt met een
% andere student?
% Zo ja, haal dan de tekst hieronder uit commentaar en pas aan.

%\paragraph{Opmerking}

% Dit voorstel is gebaseerd op het onderzoeksvoorstel dat werd geschreven in het
% kader van het vak Research Methods dat ik (vorig/dit) academiejaar heb
% uitgewerkt (met medesturent VOORNAAM NAAM als mede-auteur).
% 

\section{Inleiding}%
\label{sec:inleiding}

Het Beveiligen an beheerderstoegang tot servers vormt een crutiale uitdaging binnen moderne IT-omgevingen. Systeem- en netwerkadministratoren verbinden dagelijks aan toestellen zoals servers via protocollen zoals Secure Shell en Remote Desktop Protocol om zo de toestellen te beheren en onderhouden. Deze vormen van toegang vereisen echter verhoogde rechten, waardoor zij een aantrekkelijk doelwit vormen voor misbruik. Dit kan zowel misbruik zijn van externe aanvallers, maar evengoed door interne fouten of nalatigheid. Als beheertoegang langdurig wordt toegekend en onvoldoende wordt gelogd neemt het risico op incidenten aanzienlijk toe.

De aanleiding van dit onderzoek ontstond tijdens een interview met een IT-medewerker van Ahold Delhaize group. Dit interview vond oorspronkelijk plaats in kader van agile werken voor het vak Buisness analyse. Tijdens het interview kwam naar voren dat Ahold Delhaize een grootschalig project rond Privileged Access Managemnt had opgestart naar aanlleiding van een ernstige beveiligingsincident het jaar voordien. Volgens de informatie die de medewerker mij op dat moment kon geven bleek dat er onvoldoende zicht was op wie wanneer beheerderstoegang gebruikte. Dit verhaal inspireerde mij om dieper in te gaan op de rol van Privileged Access Management bij het beveiligen van servertoegang. Ik ging zelf checken of dit inderdaad een vaker voorkomend probleem is door te kijken bij de owasp top 10 en andere ranglijsten van veel voorkomende IT zwakten.

Uit mijn onderzoek blijkt dat dit inderdaad een veel voorkomend probleem. Niet enkel is het veel voorkomend maar ook is vaak de impact van het misbruiken van toegansrechten een zwakte die grote gevolgen kan hebben. Daarom heb ik volgende onderzoeksvraag geformuleerd: Hoe kan een Privileged Access Managment tool bijdragen aan de beveiliging van servers?

Het hoofddoel van dit onderzoek is bepalen hoe beheerderstoegang tot servers via Secure Shell (SSH) en Remote Desktop Protocol (RDP) beter kan worden beveiligd. Dit doormiddel van een Privileged Access Management-tool (PAM). Bij het bepalen van de geschikte PAM-tool zal een grote vereiste zijn dat beheerderstoegang niet langdurig en niet ongelogd mag zijn. Concreet betekend dit dat beheerders enkel tijdelijke toegang zullen krijgen. Ook zullen alle sessies gelogd moeten worden. Tot slot zullen permenente verhoogde toegangsrechten vermeden worden.

Dit onderzoek wordt uitgevoerd in de vorm van toegepast onderzoek en zal bestaan uit twee delen. Enerzijds zal er een literatuurstudie worden uitgevoerd waarin bestaande kennis, richtlijnen en martanalyes omtrend PAM worden uitgevoerd. Anderzijds zal er een proof-of-concept uitgewerkt worden waarin concreet wordt aangetoond hoe een PAM-tool de risicos omtrend beheerderstoegang via SSH en RDP kan verminderen.

De doelgroep van deze bachelorproef bestaat uit systeem-en netwerkbeheerders binnen organisaties die hun infrastructuur in eigen handen hebben.

Om een gestaafd antwoord te geven op de onderzoeksvraag zal ik ook antwoord bieden op volgende deelvragen:
\begin{itemize}
	\item Welke beveiligingsrisico’s zijn verbonden aan langdurige en ongelogde beheertoegang tot servers via SSH en RDP?
	\item Welke functionaliteiten bieden Privileged Access Management-oplossingen om tijdelijke en gelogde beheertoegang af te dwingen?
	\item In welke mate verschillen PAM-oplossingen van elkaar op het vlak van sessielogging en sessierecording?
	\item Hoe ondersteunen PAM-oplossingen Just-in-Time toegang en het principe van Zero Standing Privileges?
	\item In welke mate kan een PAM-oplossing praktisch worden toegepast om beheertoegang tot servers via SSH en RDP te beveiligen binnen een proof-of-concept?
\end{itemize}


%---------- Stand van zaken ---------------------------------------------------

\section{Literatuurstudie}%
\label{sec:literatuurstudie}

In verschillende academische en professionele literatuur wordt bevoorrechte toegang als een fundamenteel aandachtspunt beschouwd. Privileged of administrator accounts beschikken over geavanceerde rechten die je toelaten om configuraties te wijzigen, beveiligingsmechanismen te omzeilen en directe toegang te verkrijgen tot infrastructuur zoals servers. Misbruik van deze accounts wordt ook geïdentificeerd als een oorzaak met hoge impact bij beveiligingsincidenten. (oluoha et al., 2025)

Internationale richtlijnen bevestigen dit beeld. Het NIST Special Publication 800⁻53 framework benaddrukt dat onvoldoende controle van beheerdersrechten een risico vormt voor integriteit en acountability(NIST SP 800-53)

Ook binnen de OWASP Top 10 sluit deze problematiek nauw aan bij meerdere categorieën. Vooral broken Access Control en Identification and Authentication Failures. Deze worden beschreven als veelvoorkomnede oorzaken van beveiligingsincidenten. OWASP benadrukt dat het toekennen van overmatige rechten aanvallers in staat stelt om lateraal door systemen te bewegen. 

Binnen dit kader wordt Privileged Access Management geposisioneerd als een gespecialiseerde beveiligingsoplossing die gericht is op het beschermen, beheren en monitoren van verhoogde toegang. Gartner definieert PAM als technologie die beheerderstoegang controleert. Dit door het beheren van accounts, credentials en sessies die gebruikt worden voor het beheren van zowel systemen als applicaties (Gartner, 2025)

In zowel academische literatuur als in de Gartner Magic quadrant wordt er aangegeven dat klassieke toegenbeheersystemen tekortschieten. Dit omdat zij meestal geen zicht bieden op wat gebruikers met verhoogde rechten effectief uitvoeren na succesvolle authenticatie.

Een van de belangrijkste technische aandachtspunten binnen PAM is privileged session Management. Hierbij wordt beheerderstoegang tot servers centraal gefaciliteerd via gecontroleerde sessies. Systeem- en netwerkbeheerders maken hierbij gebruik van protocollen zoals SSH en RDP om systemen te beheren. De literatuur wijst erop dat onbeveiligde of niet-gelogde sessies een risico vormen. Sessies die niet gelogt worden kunnen namelijk achteraf niet getraceerd worden. Zo blijft misbruik vaak onopgemerkt (oluoha et al., 2025).

Ook dit gebrek aan zichtbaarheid sluit nauw aan bij enkele OWASP-categorieën. Het voornaamste voorbeeld hiervan is Security Logging and Monitoring Faliures. Hier is het ontbreken van voldoende logging en monitoring aangeduid als een belangrijke oorzaak van laattijdige detectie van incidenten.

Om de risico's die hierboven staan vermeld te beperken, benadrukken recente studies het belang van Just-in-Time toegang. Dit in combinatie van het vermijden van langdurige beheerdersrechten. Gartner beschrijft Just-in-Time privilege management als een kernfuncionaliteit van moderne PAM-tools. 

Deze aanpak past perfect ook binnen het principe van least privilege zoals aanbevolen door NIST. Dit verlaagt ook het risico op privilege escalation en misbruik van gecompromitteerde accounts volgens OWASP.

Verder wordt er in de literatuur sterk ingezet op loggins, auditing en sessierecording als essentiële componenten van effectieve PAM-implementaties.

Marktanalyses tonen aan dat PAM-tools zich steeds verder ontwikkelen, maar dat er aanzienlijke verschillen zijn tussen vendors. De Gartner Magic Quadrant en de Coice of the Customer-studie benadrukken dat organisaties vooral waarde hechten aan afdwingbare tijdelijke toegang, betrouwbare sessielogging en transparantie. Dit terwlijk implementatiecomplexiteit en inconsitent gedrag tussen tools maar aandachtspunten blijven (Gartner, 2025). Dit bevestigt nogmaals dat enkel het inzetten van PAM niet genoeg is. Maar dat de concrete configuratie bepalend is voor beveiligingswinst.

Hoewel er al veel bestaande literatuur en marktstudies zijn over het belang van PAM uitvoerig onderbouwen en bewijzen dat er concrete beveiligingsrisico's mee verholpen worden, blijft toegepast onderzoek beperkt. Er is weinig literatuur beschikbaar die PAM-functionaliteiten praktisch evalueert binnen concrete senario's zoals het afdwingen van tijdelijke, volledige gelogde beheerderstoegang. Dit via SSH en RDP. Deze bachelorproef zal daar zich daar specifiek over ontvermen door theoretische inzichten uit de literatuur te combineren met een proof-of-concept waarin PAM-functionaliteiten concreet worden getest en geëvalueerd in relatie tot bekende kwetsbaarheden.

%---------- Methodologie ------------------------------------------------------
\section{Methodologie}%
\label{sec:methodologie}

Deze bachelorproef kan worden gecategoriseerd als een vergelijkende studie gecombineerd met een proof of concept. Het onderzoek bestaat uit twee duidelijk afgebakende hoofddelen. In het eerste deel worden verschillende PAM-tools met elkaar vergeleken op basis van vastgestelde criteria. In het tweede deel wordt een proof-of-conceptomgeving opgezet waarin een commerciële PAM-tool wordt vergeleken met een open-source alternatief. Deze aanpak laat toe om zowel theoretische verschillen als praktische effecten van PAM-tools objectief te evalueren.

voor het beantwoorden van de onderzoeksvragen zal ik gebruik maken van volgende onderzoekstechnieken

\begin{itemize}
	\item Literatuurstudie van academische publicaties, internationale standaarden en professionele marktanalyses (o.a. Gartner, NIST, OWASP)
	\item Vergelijkende studie op basis van vooraf gedefinieerde en meetbare criteria
	\item Experiment en proof of concept, waarbij PAM-oplossingen technisch worden geïmplementeerd en getest in een gecontroleerde omgeving
\end{itemize}


Deel 1: Vergelijkende studie van PAM-vendors
Doel en scope
Het eerste onderzoeksdeel heeft als doel verschillen tussen bestaande PAM-tools systematisch in kaart te brengen. De vergelijking richt zich op marktleiders die besproken en bestudeerd worden in de Gartner Magic Quadrant en Gartner Voice of the Customer-rapporten. 
aanpak
De vergelijkende studie wordt uitgevoerd aan de hand van een gestructureerde literatuurstudie. Voor elke vendor die vergeleken zal worden, zal er onderzoek gedaan worden in bestaande literatuur zoals: officiële productdocumentatie, technische whitepapers en onafhankelijke marktanalyses.
De vergelijking wordt uitgevoerd aan de hand van volgende deelvragen:
\begin{itemize}
	\item Welke verschillen bestaan er tussen PAM-oplossingen op het vlak van authenticatie en toekenning van beheertoegang?
	\item In welke mate maken PAM-oplossingen gebruik van artificiële intelligentie voor monitoring en analyse?
	\item Welke PAM-oplossingen bieden ondersteuning voor cloud- en machine-toegang?
	\item In welke mate ondersteunen PAM-oplossingen veilig remote beheer, inclusief externe of third-party toegang?
	\item Welke verschillen bestaan er op het vlak van gebruiksgemak en implementatiecomplexiteit?
\end{itemize}
Resultaat van deel 1
het resultaat van dit deel zal bestaan uit:

\begin{itemize}
	\item Een vergelijkende tabel met PAM-oplossingen en hun functionele kenmerken
	\item Een kwalitatieve analyse van sterke en zwakke punten per vendor
	\item Een gemotiveerde selectie van één commerciële PAM-oplossing die wordt gebruikt in het proof-of-conceptgedeelte
\end{itemize}

Deel2: Proof of concept
Doel en onderzoekstechniek
Het tweede onderzoeksdeel bestaat uit een technische proof of concept. Hier zal een PAM-tool toegepast worden op een gesimuleerde bedrijfsomgeving om zo de impact van de PAM-tool aan te tonen. Dit zal gebeuren door een commerciële tool maar ook een open source tool op te zetten. Dit experiment maakt het mogelijk om theoretische verschillen uit de literatuur om te zetten in meetbare praktijkresultaten.

Opzet van testomgeving

De proof of concept omgeging wordt opgezet in een geïsoleerde testomgeving en bestaat minimaal uit:

\beign{itemize}
\item Eén of meerdere servers die beheerd worden via SSH en Remote Desktop Protocol
\item Een commerciële PAM-oplossing geselecteerd uit deel 1
\item Een open-source PAM-oplossing
\item Gesimuleerde beheerdersaccounts
\end{itemize}
in de omgeving wordt getest in welke mate beide oplossingen erin slagen om:

\begin{itemize}
\item langdurige beheerdersrechten te vermijden
\item Tijdelijke, Just-in-Time toegang af te dwingen
\item Beheerderssessies te loggen en op te nemen
\item Onbeheerde of directe toegang tot servers te blokkeren
\end{itemize}

Vergelijkingscriteria
de vergelijking gebeurt aan de hand van meetbare technische criteria zoals:
\begin{itemize}
\item Mogelijkheid tot Just-in-Time toegang
\item Aanwezigheid en deelniveau van sessielogging
\item Automatische intrekking van toegangsrechten
\item Gebruiksgemak bij dagelijkse beheerhandelingen
\end{itemize}

Resultaat van deel 2
het resultaat van dit deel bestaat uit:
\begin{itemize}
\item Een werkende proof-of-conceptomgeving
\item Screenshots en logbestanden
\item een vergelijkende evaluatie tussen de commerciële en open-source PAM-tool
\end{itemize}
%---------- Verwachte resultaten ----------------------------------------------
\section{Verwacht resultaat, conclusie}%
\label{sec:verwachte_resultaten}

Op basis van de literatuur en marktanalyses wordt verwacht dat uit de vergelijkende studie duidelijk zal blijken dat de verschillende PAM-tools significant verschillen in functionaliteiten en 
